%% =====================================================================
%%  Corrigés — Bilan 4 « Maintenant et bientôt »
%% =====================================================================
\corrigesBilan{4}

\begin{corrige}{1}
\begin{reponses}[4]
\item \rep{sitting}
\item \rep{writing}
\item \rep{swimming}
\item \rep{playing}
\end{reponses}
\end{corrige}

\begin{corrige}{2}
\begin{reponses}[2]
\item \rep{is dancing}
\item \rep{drinks}
\item \rep{are watching}
\item \rep{plays}
\item \rep{are you doing}
\item \rep{Does she work}
\end{reponses}
\pourquoi{\en{Look!} et \en{Shh!} signalent une action en cours ; \en{usually} et \en{every evening},
une habitude. Phrase 6 : une habitude, donc \en{Does…?}}
\end{corrige}

\begin{corrige}{3}
\begin{reponses}[4]
\item \rep{on}
\item \rep{in}
\item \rep{at}
\item \rep{on}
\end{reponses}
\end{corrige}

\begin{corrige}{4}
\begin{reponses}[6]
\item \rep{me}
\item \rep{him}
\item \rep{us}
\item \rep{them}
\item \rep{it}
\item \rep{her}
\end{reponses}
\end{corrige}

\begin{corrige}{5}
\begin{reponses}[2]
\item \rep{is going to visit}
\item \rep{is going to rain}
\item \rep{aren't going to fly}
\item \rep{Is she going to come}
\end{reponses}
\pourquoi{phrase 2 : on voit les nuages, c'est une prévision évidente : \en{be going to}.}
\end{corrige}

\begin{corrige}{6}
\begin{reponses}[3][french]
\item \rep{V}
\item \rep{F}
\item \rep{V}
\item \rep{F}
\item \rep{V}
\item \rep{F}
\end{reponses}
\pourquoi{Jack regarde un match de hockey à la télévision ; le rendez-vous est samedi à deux heures ; il
va faire 20\,°F, soit environ −7\,°C ! Texte de la piste page~\pageref{trans-B4.01}.}
\end{corrige}

\begin{corrige}{7}
\begin{reponses}[3][french]
\item \rep{V}
\item \rep{F}
\item \rep{V}
\item \rep{F}
\item \rep{V}
\item \rep{F}
\end{reponses}
\pourquoi{il pleut à Londres ; demain, Emma visite le British Museum (la comédie musicale, c'est samedi) ;
elle rentre le 4 janvier.}
\end{corrige}

\begin{corrige}{8}
Réponse modèle :
\begin{anglais}
\begin{tcolorbox}[otBase,colback=otGreenTint,colframe=otGreenTint,fontupper=\fontsize{10.4}{14}\selectfont]
This weekend, I'm going to visit my grandparents. On Saturday, we're going to have lunch together. On
Sunday, I'm going to play tennis with my friend Paul. It's going to be sunny!
\end{tcolorbox}
\end{anglais}
\retiens{0,5 point par phrase correcte : \en{be going to} + base verbale, préposition de temps juste.}
\end{corrige}

\begin{corrige}{9}
Réponses modèles :
\begin{reponses}[1]
\item \rep{I'm reading a book.}
\item \rep{It's cold and it's raining.}
\item \rep{It's on April 2nd.}
\item \rep{I'm going to visit my parents.}
\item \rep{I'd love to!} \textit{ou} \rep{I'm sorry, I can't. I'm busy.}
\end{reponses}
\end{corrige}
